\documentclass[a4paper,12pt,fleqn]{article}%fleqn pour aligner les equations à gauche
\usepackage{graphicx}
\usepackage[utf8]{inputenc}  
\usepackage[T1]{fontenc}
\usepackage{lscape}

\usepackage{tgheros,textcomp}%Police Arial
\renewcommand{\familydefault}{\sfdefault}

\setlength{\mathindent}{0pt}%indentation dans math = 0

\usepackage[top = 1cm, bottom = 2cm, left = 1cm, right = 1cm]{geometry}
\usepackage{titlesec}
\setcounter{secnumdepth}{4}
\usepackage{xcolor}
\usepackage{amssymb} %double flèche de réaction
\usepackage{mathrsfs}%farraday
\usepackage{xtab}%tableau sur plusieurs pages
\usepackage{esvect}%grands vecteurs

\usepackage[version=4]{mhchem} %equations chimiques

\usepackage{array}
\usepackage{chemfig}
\usepackage{multirow}
\usepackage{sectsty}
\usepackage{caption}
\usepackage{amsmath}
\usepackage{amssymb}
\usepackage[cal=boondoxo]{mathalfa} 
\usepackage{enumitem}
\usepackage{lastpage}
\usepackage{tcolorbox}
\usepackage{siunitx}
\usepackage[european, straightvoltages, RPvoltages]{circuitikz}
\usepackage{tikz}
\usetikzlibrary{shapes.geometric}
\usetikzlibrary{calc}
\usetikzlibrary{automata,positioning}
\usepackage{multicol}
\usepackage{eqnarray}
\usepackage{tabularx,colortbl}%colorer fond cellules

\sisetup{
	detect-all,
	output-decimal-marker={,},
	group-minimum-digits = 3,
	group-separator={~},
	number-unit-separator={~},
	inter-unit-product={~}
} %setup pour SIunitx

\usepackage{listings}%pour insérer le code python
\definecolor{codegreen}{rgb}{0,0.6,0}
\definecolor{codegray}{rgb}{0.5,0.5,0.5}
\definecolor{codepurple}{rgb}{0.58,0,0.82}
\definecolor{backcolour}{rgb}{0.95,0.95,0.92}

\graphicspath{ {/home/rweye/Documents/Enseignement/2022/Tspe/C16/Ressources/} }

\sectionfont{\fontsize{14}{8}\selectfont}
\renewcommand{\thesection}{\arabic{section}.}
\renewcommand{\thesubsection}{\hspace{1cm}\alph{subsection}.}

\title{\vspace{-1cm}\large 
	\begin{tabular}{|m{5cm}|>{\centering}m{6cm}|>{\centering\arraybackslash}m{6cm}|}
		\hline
		Nom : & \multirow{2}{6cm}{\centering \textbf{\textsc{DS : noyau et cortège électronique}}} &  Note et observation \\
        Prénom : & &\\
		\hline
\end{tabular}}
\date{}

\newpagestyle{mystyle}{%
	\setfoot{}{\thepage\,$|$\,\pageref{LastPage}}{}
}
\renewpagestyle{plain}{%
	\setfoot{}{\thepage\,$|$\,\pageref{LastPage}}{}
}

\pagestyle{mystyle}

\newenvironment{itemize*}%
{\begin{itemize}[label=$-$]%
		\setlength{\itemsep}{0pt}%
		\setlength{\parskip}{0pt}}%
	{\end{itemize}}

\newenvironment{itemize**}%
{\begin{itemize}[label=]%
		\setlength{\itemsep}{0pt}%
		\setlength{\parskip}{0pt}}%
	{\end{itemize}}

\renewcommand{\arraystretch}{1.5}


\usepackage{xparse}
\ExplSyntaxOn
\NewExpandableDocumentCommand \randomint { m }
{ \int_rand:nn { #1 - #1 } { #1 + #1 } }
\ExplSyntaxOff

\usepackage{arydshln}%pour des ligne tableau en pointillé

\begin{document}
\maketitle
\vspace{-2.5cm}
\noindent \begin{minipage}[t]{.95\linewidth}
\section*{Exercice 1 : Twins from two mothers}
On considère deux noyaux atomiques le bismuth 209 ($^{209}_{82}Bi$) et le polonium 209 ($^{209}_{84}Po$).
\begin{enumerate}
    \item Donner la composition des deux noyaux et les comparer.
    \item Calculer la masse des deux noyaux et les comparer.
    \item Donner dans chacun des cas le nombre d'électrons présent pour les atomes associés.
    \item S'agit-il d'isotopes ? Justifier.
\end{enumerate}
\textbf{Données : }
\begin{itemize*}
    \item $m_{nucleon}$ = \SI{1,67E-27}{kg}
    \item $m_{electron}$ = \SI{9,11E-31}{kg}
\end{itemize*}


\section*{Exercice 2 : Configuration électronique}
On considère les éléments suivants : $_{13}Al$, $_{12}Mg$, $_{15}P$.
\begin{enumerate}
    \item Écrire la configuration électronique de ces trois atomes.
    \item Justifier le positionnement de ces éléments dans le tableau périodique et compléter le tableau ci-après.
    \item Donner la définition d'un ion monoatomique.
    \item Déterminer en justifiant, quels ions monoatomiques ces trois atomes vont former.
    \item Donner la formule ionique de la fluorine, formée d'ions calcium et d'ions fluorure.
\end{enumerate}
\end{minipage}
\fbox{\begin{minipage}[t]{.05\linewidth}
\flushright
\vspace{3.5em}
    /1\vspace{1em}
    
    /2\vspace{1em}
    
    /1\vspace{1em}
    
    /1\vspace{5cm}
    
/1,5\vspace{1em}
    
/2\vspace{1em}
    
/1\vspace{1em}
    
/1,5\vspace{1em}
    
/1\vspace{1em}
\end{minipage}}

\begin{table}[h!]
\small \centering 
\caption{Tableau périodique simplifié}
\begin{tabular}{|c|>{\centering}m{1.8cm}|>{\centering}m{1.5cm}|c:c|>{\centering}m{1.3cm}|>{\centering}m{1.3cm}|>{\centering}m{1.3cm}|>{\centering}m{1.5cm}|>{\centering}m{1.3cm}|c|}
		\cline{1-3}\cline{6-11} & 1 & 2 & & & 3 & 4 & 5 & 6 & 7 & 8\\
		\cline{1-3}\cline{6-11} 1 & \textbf{\ce{^1_1He}}& & & & & & & & & \textbf{\ce{^4_2He}}\\
		&Hydrogène& & & & & & & & & Hélium\\
		
		\cline{1-3}\cline{6-11} 2 & \textbf{\ce{^7_3Li}}& \textbf{\ce{^9_4Be}} & & & \textbf{\ce{^11_5B}}&\textbf{\ce{^12_6C}} &\textbf{\ce{^14_7N}} & \textbf{\ce{^16_8O}}&\textbf{\ce{^19_9F}} & \textbf{\ce{^20_10Ne}}\\
		&Lithium & Béryllium & & & Bore & Carbone & Azote & Oxygène & Fluor & Néon\\
		
		\cline{1-3}\cline{6-11} 3 & \textbf{\ce{^23_11Na}}&  & & & &\textbf{\ce{^28_14Si}} & & \textbf{\ce{^32_16S}}&\textbf{\ce{^25_17Cl}} & \textbf{\ce{^40_18Ar}}\\
		&Sodium &  & & &  & Silicium &  & Soufre & Chlore & Argon\\
		\cline{1-3}\cline{6-11}		
	\end{tabular}
\end{table}


\end{document}